\section{The scanner estimate with four norm comparisons}
\label{sec:al_scanner_source}

The scanner estimate of \cite[Proposition~9.2]{OpenAI2026AreaLaw} selects
one charge round and one interpolation parameter at which both the density
and energy bounds and all four norm comparisons hold. The coefficients in
the output shell and density bounds may be enlarged after the input entropy
coefficient is fixed, as recorded in~\cite{gap:arealaw2d_scanner_inputs}.

% Source: 08-scanner.tex, prop:scanner, lines 349--397; scales and geometry,
% lines 21--59; 07-comparators.tex, prop:comparators, lines 125--214.
% Revision: adc7f1241b42e322a6451854ab7e4b4c146bf78a.
\begin{definition}[Statement of the scanner estimate]
    \label{def:al_scanner_source_statement}
    \lean{TNLean.PEPS.AreaLaw.Scan.CollarScan.ScannerHypotheses,
        TNLean.PEPS.AreaLaw.Scan.CollarScan.typicalWindow,
        TNLean.PEPS.AreaLaw.Scan.CollarScan.scanPrevector,
        TNLean.PEPS.AreaLaw.Scan.CollarScan.ScannerStatement}
    \leanok
    \uses{thm:al_scan_truncation_inputs,thm:arealaw-supplied-target-norm-bounds}
    Fix exponents $0<\ell<1$, $0<\kappa<\mu<1-\ell$, $0<e<1$ with
    $\kappa\le e$, and $\nu>0$; an interaction range $R$, a term bound $J$,
    a gap $\Delta$ and a target-size coefficient $C_T$. Put
    $g=\Delta/\max\{1,J(\|f_{1,\Delta/2}\|_1+1)\}$. At scale $n$, put
    \begin{align}
        L & =\lfloor n^{1-\ell}\rfloor,
          & m                           & =\lfloor n^\mu\rfloor,
          & K                           & =\lfloor L/(8m)\rfloor,\notag \\
        D & =\lceil n^\kappa\rceil,
          & \epsilon                    & =n^{-\nu},
          & a                           & =W/K.\notag
    \end{align}
    The source hypotheses at scale $n$ are the following. The unit vector
    $\Omega$ is a ground vector, with gap $\Delta$, of a local Hamiltonian
    on $\Lambda$ with range $R$ and term bound $J$. The scan uses the
    induced lattice graph, its depth is the ambient sup-norm distance to
    $T$, and every designated support contains its anchor. The collar scan has $n$
    slots per depth row, the displayed scales, truncation radius
    $r_0=\lceil C_{\rm tr}(\log n)^2\rceil$ and charge-list constant $C_1$.
    Its finite target $T$ obeys $|T|\le C_Tn^2$ and
    $|T_j\setminus T_{j-1}|\le n$ for $1\le j\le L$. The ambient distance
    from $T$ to every endpoint of a boundary edge of $A$ exceeds
    $2L+10r_0$. The target $X=A\cap T$ is nonempty. With
    $Q_j=A\cap(T_j\setminus T)$, the original vector satisfies
    \begin{align}
        S_\Omega(X)   & \le C_en^{1+e},\notag \\
        S_\Omega(Q_j) & \le C_e(nL^e+n)\notag
    \end{align}
    for every integer $0\le j\le L$.

    The scanner statement asserts that there is $C_{\rm tr}>0$ such that,
    for every integer $q\ge2$ and every real $C_e$, there are constants
    $C_1>0$, $c_a>0$, $C>0$, $C_\delta>0$, $C_b,C_{l,b},C_l,C_{l,\delta}$
    and a size threshold with the following property. For every larger $n$ and every
    finite physical system satisfying the source hypotheses, the truncated
    Hamiltonian has a unit ground vector $\widetilde\Omega$ of energy
    $0\le e_t\le n^{-1000}$ and gap $g/2$. Its target typical set $E$ at
    width $n^{3/5}$ is nonempty. Write $d_*=|E|$ and let $z$ be its mass;
    then $0<1-n^{-100}\le z\le1$, in particular $z>0$, for this same
    $\widetilde\Omega$ and $E$.
    There is one label sequence $\lambda_k$ such that
    $\mathrm{pre}_k=\pi_R^{\lambda_k}
        (\widetilde\Omega\otimes\mathrm{Bell}_{CR})^{\otimes k}$
    is nonzero for $k>0$, with both auxiliary dimensions equal to $d_*$.
    These choices precede the total weight $W$.

    For every $W\ge1$ with $ar_0^2\le c_a$, there are comparator parameters
    \begin{align}
        B_{\rm sh}     & \le C_b(nL^e+n),\notag             \\
        B_{\rm exc}    & \le C_bnD,\notag                   \\
        1\le\mathcal B & \le C_bnD(\log n)^{C_{l,b}}.\notag
    \end{align}
    Every completed and pre-charge status cut has labelled budget at most
    $\mathcal B$. At each physical region, the marginal moment estimate
    for $\widetilde\Omega$ uses that region's own labelled cut budget. For all sufficiently large replica counts $k$, the final charge
    at interpolation parameter $1$ satisfies
    \begin{align}
        -\frac{\log N_{\rm end,k}^2}{k}
         & \le2\log d_*-\log z
        +W(B_{\rm sh}/z+B_{\rm exc})+\rho_{\rm end,k},\notag
    \end{align}
    where $\rho_{\rm end,k}\to0$.

    Set
    \begin{align}
        \delta_n & =C_\delta\epsilon^{-1}
        \left(n^{e-\mu}+a^{1/4}(\log n)^{C_{l,\delta}}\right).\notag
    \end{align}
    There is $\rho_k\to0$ such that, for every sufficiently large $k$, one
    charge round and one $p_k\in[\epsilon/2,\epsilon]$ satisfy
    \begin{align}
        \frac{\mathcal Q_k(p_k)}{KnD} & \le\delta_n+\rho_k,\notag        \\
        E_{{\rm def},k}               & \le Cn^\ell D^2W^2(\log n)^{C_l}
        \left(\delta_n^{1/8}+\epsilon+n^{-100}\right)
        +Cn^{-1000}+\rho_k.
        \label{eq:al_scanner_source_energy}
    \end{align}
    Here $v_k$ is the normalized filtered vector at this charge and
    parameter, and
    $E_{{\rm def},k}=(\langle v_k,\overline{\widetilde H}_k v_k\rangle-e_t)/(g/2)$.
    The filtered norm is positive and $E_{{\rm def},k}\ge0$.

    Use the normalized Schmidt truncation $\psi$, its auxiliary entropy
    $S^\prime$, and the terminal-region functions and charge-tree weights
    of Theorem~\ref{thm:arealaw-supplied-target-norm-bounds}, for these same
    physical data. At the selected point, the rough upper comparison is
    \begin{align}
        -\frac{\log N_k^2}{k} & \le2\log d_*-\log z
        +W(B_{\rm sh}/z+B_{\rm exc})+R_k^{\rm RU},\notag
    \end{align}
    and the rough lower comparison is
    (\ref{eq:arealaw-supplied-rough-lower}), for every $0<\tau\le1/2$.
    With $C_0=512\mathrm e(2/g)$ and
    $c_0=[32\sqrt{1+2/g}]^{-1}$, the additional condition
    $a\sqrt{\mathcal B}\le\min\{c_0,1\}/10$ gives both sharp comparisons
    (\ref{eq:arealaw-supplied-sharp-bounds}). Their state, typical set,
    labels, charge and interpolation parameter are the same. The norm
    remainders tend to zero with the finite system and $W$ fixed and are
    uniform in the selected charge and parameter.

    The input coefficient $C_e$ is unchanged. The output shell coefficient
    may dominate $\max\{C_e,1\}+1$, and the density coefficient
    $C_\delta$ may be positive when $C_e=0$. The assertion with both output
    coefficients numerically identical to $C_e$ is not claimed;
    see~\cite{gap:arealaw2d_scanner_inputs}.
\end{definition}

\begin{theorem}[The scanner estimate]
    \label{thm:al_scanner_source}
    \lean{TNLean.PEPS.AreaLaw.Scan.CollarScan.scannerStatement}
    \leanok
    \uses{def:al_scanner_source_statement}
    The scanner statement in Definition~\ref{def:al_scanner_source_statement}
    holds for $J\ge0$, $\Delta>0$ and $C_T\ge0$. It gives the conclusion of
    \cite[Proposition~9.2]{OpenAI2026AreaLaw} with all four norm comparisons
    and the coefficient convention in~\cite{gap:arealaw2d_scanner_inputs}.
\end{theorem}

\begin{proof}
    \leanok
    \uses{thm:arealaw-source-selected-energy,thm:arealaw-supplied-target-norm-bounds,
        thm:al_scan_common_status_budget,thm:al_positive_prefix_entropy_transfer,
        thm:al_scan_actual_rough_upper,thm:al_augmented_prevector_norm,thm:al_lattice_counts}
    Choose the truncation coefficient from
    Theorem~\ref{thm:arealaw-source-selected-energy}, before fixing $q$.
    The anchor of a support of range $R$ has multiplicity at most
    $\mu_*=2^{1+2R(R+1)-1}$, so take $C_1=3\mu_*+1$.
    Apply that theorem with entropy coefficient $\max\{C_e,1\}$.
    It supplies one truncated vector, typical set and label sequence and,
    for each admissible $W$, one selected charge and parameter with the
    density, energy and rough upper bounds.
    Retain its typical-window estimate $0<1-n^{-100}\le z\le1$ for
    that same vector and set, as in \texttt{08-scanner.tex}, lines~411--414.

    Apply Theorem~\ref{thm:arealaw-supplied-target-norm-bounds} to those
    same data. Its rough lower bound is uniform over charge rounds and
    parameters. The common cut estimate and the norm estimate give
    coefficients $C_{B,\rm cut}$ and $C_{B,\rm norm}$. Take
    \begin{align}
        C_b        & =\max\{C_{B,\rm cut},C_{B,\rm norm},
        \max\{C_e,1\}+1,1+3\log q\},\notag                \\
        \mathcal B & =C_bnD(\log n)^{12}.\notag
    \end{align}
    Enlarging the marginal parameter preserves the sharp upper and lower
    estimates and strengthens their smallness antecedent.

    Entropy continuity on each whole-row positive-depth prefix gives
    $B_{\rm sh}=\max\{C_e,1\}(nL^e+n)+n^{-247}$; take
    $B_{\rm exc}=1+3nD\log q$. Since $n\ge2$ and $D\ge1$, these obey the
    required bounds. The uniform rough upper estimate at the final charge
    gives the last-round assertion with its own vanishing remainder.
    Intersect the eventual selected-point and sharp-comparison estimates,
    retaining the point already chosen by the energy theorem. The sum of
    its nonnegative density and energy remainders serves both selected
    inequalities. This leaves every physical state, label and selected
    point unchanged.
\end{proof}
